\section{Disucussion}\label{sec:threats}
\subsection{Limitations}
We identify three limitations of \name: 1) the incapacity to deal with bugs related to program logic; 2) the prerequisites for multi-source data collection; 3) the requirement of annotated data for training.

As \name is an entirely data-driven approach targeting the scope of reliability management, it is only applicable to troubleshooting anomalies manifested in the involved data, so logical bugs out of our scope and silent issues that do not incur abnormal patterns in observed data can not be detected or located.

Moreover, \name is basically well-suited for all microservices where anomalies can be reflected in the involved three types of data we employ. However, some practical systems may lack the ability to collect the three types of data. Though the low-coupled nature of the modal-wise learning module allows the absence of some source of data, it is better to provide all data types to fully leverage \name.
Since we apply standard open-source monitoring toolkits and these off-the-shelf toolkits can be directly instrumented, enabling microservices with the data collection ability is not difficult.

In addition, the supervised nature of \name requires a large amount of labeled training data, which may be time-consuming in the real world. 
Nevertheless, our approach outperforms compared with unsupervised approaches by a large margin, indicating that in practice, unsupervised methods may be difficult to use because the accuracy rate is not up to the required level, especially considering that realistic microservices systems are much larger and more complex.
A common solution in companies is to use an unsupervised model to generate coarse-grained pseudo-labels. Afterward, experienced engineers manually review the labels with lower confidence.
The hybrid-generated labels are used for training the supervised model, and eventually, the supervised approach performs the troubleshooting work.
Hence, \name will still play an important role in practice and fulfill its potential.

\subsection{Threat to Validity}
\subsubsection{Internal Threat}
The main internal threat lies in the correctness of baseline implementation.
We reproduce the baselines based on our understanding of their papers since most baselines, except DyCause and TraceAnomaly, have not released codes, but the understanding may not be accurate.
To mitigate the threat, we carefully follow the original papers and refer to the baseline implementation released by~\cite{DyCause}. 

\subsubsection{External Threat}
The external threats concern the generalizability of our experimental results. We evaluate our approach on two simulated datasets since there is no publicly available dataset containing multi-modal data.
It is yet unknown whether the performance of \name can be generalized across other datasets.
We alleviate this threat from two aspects.
First, the benchmark microservice systems are widely used in existing comparable studies, and the injected faults are also typical and broadly applied in previous studies~\cite{MicroRank, AutoMAP, DyCause}, thereby supporting the representativeness of the datasets.
Second, our approach is request- and fault-agnostic, so an anomaly incurred by a fault beyond our injections can also be discovered if it causes abnormalities in the observations.
